\documentclass[11pt,a4paper]{article}
\usepackage[a4paper,top=22mm,bottom=22mm,left=23mm,right=23mm]{geometry}
\usepackage{amsmath,amssymb,xcolor,hyperref}
\definecolor{ink}{RGB}{20,63,69}
\hypersetup{colorlinks=true,urlcolor=ink,linkcolor=ink,pdfauthor={Lee Crossley},pdftitle={Sign changes in many residue classes at every prime modulus},pdfsubject={Research note on integrated prime-counting errors}}
\setlength{\parindent}{0pt}
\setlength{\parskip}{5pt plus .5pt minus .5pt}
\setlength{\abovedisplayskip}{6pt plus 1pt minus 1pt}
\setlength{\belowdisplayskip}{6pt plus 1pt minus 1pt}
\setlength{\emergencystretch}{1.8em}
\newcommand{\li}{\operatorname{li}}
\newcommand{\FP}{\operatorname{FP}}
\newcommand{\kb}{\overline{\chi}}
\newcommand{\head}[1]{\section*{\color{ink}#1}}
\raggedbottom
\begin{document}
\begin{center}
{\Large\bfseries Sign changes in many residue classes\\[2pt]at every prime modulus\par}
\vspace{6pt}
{Lee Crossley \quad$\cdot$\quad Johnston's Problem 23 \quad$\cdot$\quad 21 November 2026}
\end{center}

Let \(q\) be prime, \(h=q-1\), and let \(\pi(t;q,a)\) count primes \(p\le t\) in
the nonzero residue class \(a\bmod q\). Write
\[
 E_{q,a}(x)=\frac{\log x}{x^{3/2}}\int_2^x\bigl(h\pi(t;q,a)-\li(t)\bigr)\,dt .
\]
The question is whether \(E_{q,a}\) changes sign infinitely often in every class at every
prime \(q\ge11\). We give a partial answer for unbounded \(q\).

\textbf{Theorem.} Assume the Riemann hypothesis for the primitive Dirichlet functions
inducing the characters involved, including \(\zeta\). There are absolute
\(\delta,\eta>0\) such that, for every prime \(q\ge11\), at least \(\delta(q-1)\)
distinct classes satisfy
\begin{equation}
 \liminf_{x\to\infty}E_{q,a}(x)<-\eta\sqrt{q\log q},
 \qquad \limsup_{x\to\infty}E_{q,a}(x)>\eta\sqrt{q\log q}.
\end{equation}
For all sufficiently large primes the assertion holds separately for a positive
fraction of the squares and of the nonsquares. More generally, for fixed \(d\mid q-1\),
it holds in every coset of the subgroup of index \(d\) as \(q\) increases, with
constants depending on \(d\). Across all classes at a prime \(q\), at most
\(O((\log q)^2)\) can fail to have a strictly negative lower limit.

Under the same hypothesis a particular class also has both signs for every prime \(q\ge11\). For \(q\le10\,009\) this also follows unconditionally from the finite certificate. Take \(a=2\) if
\(q\equiv1\bmod3\), and \(a=4\) if \(q\equiv2\bmod3\). For \(q\ge10\,000\)
its upper limit exceeds \(q/10\) and its lower limit is below \(-\sqrt q/60\).
The same congruence gives a growing list of explicit powers of two on page 6.

The argument allows central zeros, coincident zeros and rational relations among
their heights. The existence constants in (1) are not evaluated. Alpöge's announced Dirichlet
GRH theorem supplies this input, and Sutherland records independent formal checks.
The present argument is a written proof awaiting independent review.
\href{https://alpo.ge/remark.pdf}{Announcement.} \href{https://math.mit.edu/~drew/grh/}{Independent formal checks.}

\head{The profile and its recurrence}
Let \(C_a\) count the square roots of \(a\bmod q\), put \(B_a=2C_a/3\), and use
\(\chi^*\) for the primitive inducing character. For \(\rho=\tfrac12+i\gamma\)
ranging over all nontrivial zeros of its function with multiplicity, set
\begin{equation}
 F_a(u)=-B_a-\sum_{\chi\bmod q}\kb(a)\sum_{\rho}
           \frac{e^{i\gamma u}}{\rho(\rho+1)} .
\end{equation}
The first Riesz mean explicit formula, removal of prime powers and partial summation
give \(E_{q,a}(x)=F_a(\log x)+o_q(1)\).
The zero count gives absolute uniform convergence in (2). Every strict value of
this continuous almost periodic function recurs to arbitrary precision at arbitrarily
large times. Thus a pair of strict profile values yields arbitrarily large inputs
of both signs. After decreasing the margins slightly, the two sets of inputs each
have positive lower logarithmic density for fixed \(q,a\). The fixed-modulus
explicit formula and density argument are in
\href{https://github.com/leecrossley/integrated-prime-errors-mod-11/blob/v2.0.0/range/PROOF.md}{the previous paper, \S1}.

\newpage
\head{Two arithmetic evaluations of the same profile}
Use the representatives \(1\le a,b<q\) with \(ab\equiv1\bmod q\), and suppose
\(a\ne\pm1\bmod q\). Write \(\Lambda\) for the von Mangoldt function and define
\[
 \begin{aligned}
 D_a(s)&=-\sum_{\chi\bmod q}\kb(a)\frac{L'}{L}(s,\chi^*)
             =\sum_{n\ge1}\frac{w_a(n)}{n^s}\quad(\Re s>1),\\
 w_a(n)&=h\,1_{n\equiv a\bmod q}\Lambda(n)+1_{n=q^j,\;j\ge1}\log q,\\
 c_b&=\FP_{s=1}D_b(s),\qquad d_b=D_b(2),\qquad
 M_a(y)=\sum_{n\le y}w_a(n)(y-n).
 \end{aligned}
\]
The principal character contributes the primitive factor \(\zeta\).
The weighted sum over each parity vanishes because \(a\ne\pm1\).
In the primitive functional equations this removes the gamma factors and leaves
the weighted conductor sum \(-\log q\). Consequently
\[
                 D_a(s)+D_b(1-s)=-\log q.
\]
For \(y\ge1\) the two identities we use are
\begin{align}
 \sqrt y\bigl(F_a(\log y)+B_a\bigr)
    &=\frac{M_a(y)}y-\frac y2+1+\log(q/y)+c_b-\frac{\log q+d_b}{y},\\
 \frac{F_a(-\log y)+B_a}{\sqrt y}
    &=1+\log y+c_b-\sum_{n\le y}\frac{w_b(n)}n
                    -y\sum_{n>y}\frac{w_b(n)}{n^2}-\frac1{2y}.
\end{align}
The tail in (4) equals \(d_b-\sum_{n\le y}w_b(n)/n^2\).
The forward formula specializes to the published
\href{https://github.com/leecrossley/integrated-prime-errors-mod-11/blob/v2.0.0/PROOF.md}{modulus-11 computation, \S6}.

For (3), shift the inverse Mellin integral of \(D_a(s)y^{s+1}/[s(s+1)]\) left.
The pole at 1 contributes \(y^2/2\). At 0 use
\(D_a(s)=1/s-\log q-c_b+O(s)\), giving \(y(\log y-1-\log q-c_b)\).
At \(-1\) the residue is \(-D_a(-1)=\log q+d_b\). The nontrivial
zero terms give \(y^{3/2}(F_a(\log y)+B_a)\).

For (4) use the kernel \(1/[(s-1)(s-2)]\) on a line to the right of 2.
Mellin inversion for \(D_b(s)y^s\) times this kernel gives
\(\sum_{n\le y}w_b(n)((y/n)^2-y/n)\).
The residues at \(2,1,0\) are \(d_b y^2,-y(1+\log y+c_b),1/2\).
Reflection \(\rho\mapsto1-\rho\) interchanges the character and its conjugate,
so the zero contribution is \(y^{1/2}(F_a(-\log y)+B_a)\).
The weighted functional equation removes all other trivial poles in both shifts.
For fixed \(q\), take horizontal heights separated from ordinates by
\(\gg_q1/\log T\), using the local zero count. The usual logarithmic derivative
bound \(O_q(\log^2 T)\) on each fixed-width horizontal segment and the quadratic
denominators make those integrals tend to zero. On a far left vertical line the
reflected series is bounded uniformly. Its integral tends to zero for \(y>1\).
Continuity gives both identities at \(y=1\).

\newpage
\head{Uniform bounds and negative excursions}
Uniformly for \(1\le b<q\), under the same critical-line assertion,
\begin{equation}
 c_b=\frac{h\Lambda(b)}{b}+O(\log^2q),\qquad
 d_b=\frac{h\Lambda(b)}{b^2}+O\!\left(\frac{\log q}{q}\right).
\end{equation}
To see the first estimate put \(Y=q^2\). Fiorilli and Martin's
\href{https://arxiv.org/pdf/0912.4908#page=31}{Proposition 3.10}
approximates each \(-L'(1,\chi)/L(1,\chi)\) by the exponentially smoothed series
with individual error \(O(\log q/q)\). The principal finite part is \(-\gamma_E\).
Orthogonality then gives
\[
 c_b=-\gamma_E+h\sum_{n\equiv b(q)}\frac{\Lambda(n)e^{-n/Y}}n
                -\sum_{q\nmid n}\frac{\Lambda(n)e^{-n/Y}}n+O(\log q).
\]
Their Proposition 3.11 gives \(\Lambda(b)/b+O(\log^2q/q)\) for the progression
sum. The unrestricted smoothed sum is \(O(\log q)\) by the elementary
bound \(\sum_{n\le x}\Lambda(n)/n\ll\log(2x)\) and partial summation.
For \(d_b\), the remaining progression terms are at most
\(h\sum_{j\ge1}\log((j+1)q)/(jq)^2\), and the ramified contribution is
\(\log q/(q^2-1)\). This proves (5).

At \(y=a\) each summand of \(M_a(y)\) is zero. Formula (3) and \(d_b\ge0\)
imply, for an absolute \(K\) and \(a\ne\pm1\),
\[
 \sqrt a\,F_a(\log a)\le-\frac a2+\frac{h\Lambda(b)}b+K\log^2q .
\]
This is strictly negative if \(a>4K\log^2q\) and \(ab>4h\Lambda(b)\).
The first condition has \(O(\log^2q)\) exceptions. Failure of the second
forces \(b\) to be a prime power dividing \(kq+1=ab\) with
\(1\le k\le4\log q\). Each such integer has at most
\(\log(kq+1)/\log 2\) prime-power divisors greater than one, since their number is
the sum of the exponents of its prime factors. As \(b\) determines \(a\), this
gives another \(O(\log^2q)\) exceptions. The two omitted classes fit the bound.

\head{A single prime gives a positive excursion}
Set \(X=h\log q/8\). If \(\sqrt q\le p\le X\) is prime and \(p\equiv a\bmod q\),
then \(M_a(2p)/(2p)\ge h\log p/2\). By (5),
\(c_b-d_b/y\ge-O(\log^2q)\) for \(y\ge1\). Substitute \(y=2p\) in (3) to obtain
\[
 \sqrt{2p}\,F_a(\log(2p))
  \ge h\log p/2-p-O\!\left(\sqrt p+\log^2q\right)
  \ge h\log q/8-O\!\left(\sqrt{q\log q}+\log^2q\right).
\]
Here \(0\le B_a\le4/3\) and \(|\log(q/(2p))|\ll\log q\).
Thus \(F_a(\log(2p))\gg\sqrt{q\log q}\) for every sufficiently large \(q\)
in the indicated range, with an absolute implied constant.

\newpage
\head{Counting classes with both signs}
At \(X=h\log q/8\), write \(N_a=\#\{p\le X:p\equiv a\bmod q\}\).
Erd\H{o}s's \href{https://www.renyi.hu/~p_erdos/1949-05.pdf}{Theorem 2}
gives at least \(\delta_0 h\) classes with \(N_a>0\) for some absolute \(\delta_0>0\).
For the refinement, let \(H\) be the subgroup of fixed index \(d\), and \(C\) one
of its cosets. Equation (9) of the same paper gives
\(\sum_a N_a(N_a-1)\ll h\). Since \(\sum_aN_a\le\pi(X)\ll h\), the second moment
\(\sum_aN_a^2\ll h\) as well.

Character orthogonality on the quotient by \(H\) gives
\[
 \sum_{a\in C}N_a=\frac{\pi(X)}{d}+O\!\left(\sqrt X\log^2(qX)+1\right)
                 \sim\frac{h}{8d}.
\]
Indeed there are \(d\) characters on the quotient. For each nonprincipal one
the GRH estimate \(\sum_{n\le x}\Lambda(n)\chi(n)\ll\sqrt{x}\log^2(qx)\)
and prime-power removal followed by partial summation give the displayed error.
The principal term omits the prime \(q\). See Montgomery and Vaughan,
\emph{Multiplicative Number Theory I}, Theorem 13.7. Cauchy--Schwarz now gives
\(\#\{a\in C:N_a>0\}\ge(\sum_{a\in C}N_a)^2/\sum_aN_a^2\ge\delta_d h\)
for an absolute-in-\(q,C\) constant \(\delta_d>0\) for each fixed \(d\).

Choose fixed \(0<\varepsilon<\min(1/16,\delta_d/8)\) and set
\(Y=\varepsilon h\log q\). The number of classes containing any prime
power \(\le Y\) is at most \(\pi(Y)+O(\sqrt Y\log Y)\sim\varepsilon h\), by the
ordinary prime number theorem. Discard these classes and \(a=\pm1\). At least
\(\delta_d h/2\) classes in \(C\) remain for large \(q\), each with a prime
\(Y<p\le X\). Such a prime exceeds \(\sqrt q\), so the preceding page supplies
the positive excursion.

For these same classes the only possible contribution to \(M_a(Y)/Y\) comes from
the powers of \(q\), of total weight at most \(\log Y\). From (3), (5),
\(\Lambda(b)/b\le1\) and \(d_b\ge0\), we obtain
\[
 \sqrt{Y}\,F_a(\log Y)\le -Y/2+O(q+\log^2q+\log Y)\le -Y/3
\]
for sufficiently large \(q\), since \(Y/q\to\infty\). This supplies a negative
excursion of size \(\gg_d\sqrt{q\log q}\). Recurrence proves (1), with strict
margins obtained by decreasing the constants. Take \(d=2\) for squares and
nonsquares, or \(d=1\) for the entire group.

For the unqualified group the class \(1\bmod q\) has a strictly negative lower
limit and a strictly positive upper limit at every prime \(q\ge11\), by the
\href{https://github.com/leecrossley/integrated-prime-errors-mod-11/blob/v2.0.0/range/PROOF.md}{previous class-one theorem}.
Decreasing \(\delta,\eta\) over the finitely many intervening primes extends (1)
from sufficiently large \(q\) to every \(q\ge11\). The resulting constants are
existence bounds. There is no numerical claim about their size.

\newpage
\head{An explicit class at every prime}
For \(q\equiv1\bmod3\) choose \(a=2\), and for \(q\equiv2\bmod3\) choose \(a=4\).
In both cases \(\Lambda(a)=\log2\) and \(3\mid q+a\). For \(q\ge11\) the
only possible progression terms up to \(3q\) are \(a,q+a,2q+a\).
The middle term contributes weight \(\log3\) only if it is a power of 3.
The last term is even and can contribute only weight \(\log2\).

These last two terms cannot both contribute. For \(a=4\), the last is
\(2(q+2)\) with \(q+2>1\) odd. For \(a=2\), simultaneous contributions would
give \(q+1=2^m\), \(q+2=3^k\). Here \(m\ge3\), so \(3^k\equiv1\bmod8\)
forces \(k\) even. The two positive even factors \(3^{k/2}-1,3^{k/2}+1\) of
\(2^m\) are powers of 2 differing by 2. Only \(2,4\) do so, giving \(q=7\).
Consequently
\[
 \frac{M_a(3q)}{3q}\le h\left(\log2+\frac23\log3\right)+\frac23\log q,
 \qquad \frac{M_a(2a)}{2a}\ge\frac{h\log2}{2}.
\]
The common contribution at \(n=q\) accounts for the last log in the first bound.

Here are explicit bounds for \(q\ge10\,000\). Put \(l=\log q,\ U=h\Lambda(b)/b\).
The two Fiorilli--Martin bounds used in (5) give an error at most
\[
 2l^2+3.935l+(14.27l+16.25)+(16.1l+17.83)/\sqrt q
 <2l^2+19l+17
\]
in \(c_b=U-\gamma_E-\sum_{q\nmid n}\Lambda(n)e^{-n/q^2}/n+\text{error}\).
The subtracted sum is at most \(2\log(1+q^2)<4l+1\). Indeed
\(\log(m!)=\sum_{n\le m}\Lambda(n)\lfloor m/n\rfloor\) implies
\(\sum_{n\le m}\Lambda(n)/n\le2\log m\) for \(m\ge1\).
Average using a geometric random integer \(m\ge0\) with
\(\Pr(m\ge n)=e^{-n/q^2}\). By concavity,
\(2\,\mathbb E\log(1+m)\le2\log(1+q^2)\).
Also \(0\le d_b-U/b\le(2\log(2q)+2)/q+l/(q^2-1)<1\). Since \(ab=1+kq\)
with \(k\ge1\), \(b>q/a\) and \(U\le4l\).

With \(D=3/2-\log2-\tfrac23\log3\), formula (3) now gives
\[
 \begin{aligned}
 \sqrt{2a}\,F_a(\log(2a))&\ge h\log2/2-2l^2-24l-28,\\
 \sqrt{3q}\,F_a(\log(3q))&\le-Dq+2l^2+24l+17 .
 \end{aligned}
\]
For the first line use \(4\le2a\le8\), \(B_a\le4/3\), and
\(c_b-d_b/(2a)\ge-2l^2-23l-20\).
At \(q=10\,000\) these bounds give \(F_a(\log(2a))/q>0.1077\) and
\(-F_a(\log(3q))/\sqrt q>0.0194\). Both bounds improve as \(q\) increases,
since \((2l^2+24l+17)/q\) and \((2l^2+24l+28)/q\) decrease.
Recurrence gives the claimed margins \(1/10\) and \(1/60\). The published
unconditional range covers the smaller primes. A separate exact sieve checked
the prime-power exclusion at all 148,929 primes from 11 to two million.

\newpage
\head{Further powers of two and the limit of a single prime}
Let \(q\ge10\,000\), and take any \(a=2^j\), \(j\ge1\), satisfying
\[
                     3\mid q+a,\qquad a\le\frac{q}{50\log q}.
\]
Their representatives are smaller than \(q\), so these classes are distinct.
The congruence selects every other exponent, giving
\((1/2+o(1))\log_2 q\) such classes as \(q\to\infty\).
Here \(\log_2 q=\log q/\log2\). Each satisfies
\[
     \limsup E_{q,a}>\frac{q}{4\sqrt{2a}},
     \qquad \liminf E_{q,a}<-\frac{\sqrt q}{120}.
\]
For \(j\ge2\), the integer \(2q+a=2(q+2^{j-1})\) has an odd factor greater than one.
For \(j=1\) the preceding page treats the possible power of 2. Thus exactly the
same upper bound for \(M_a(3q)/(3q)\) applies, and
\(M_a(2a)/(2a)\ge h\log2/2\) still holds.
Since \(b>q/a\), the bound \(U\le al\le q/50\) suffices.
For \(y=2a\) we have \(B_a\sqrt{2a}\le4a/3\) and \(\log(q/(2a))>0\).
Substitution into (3) gives
\[
\begin{aligned}
 \frac{\sqrt{2a}\,F_a(\log(2a))}{q}
  &\ge\left(1-\frac1q\right)\frac{\log2}{2}
        -\frac{2l^2+24l+20}{q}-\frac7{150l},\\
 \frac{-F_a(\log(3q))}{\sqrt q}
  &\ge \frac1{\sqrt3}\left(D-\frac1{50}
                     -\frac{2l^2+20l+17}{q}\right).
\end{aligned}
\]
The right sides increase for \(q\ge10\,000\) and exceed \(0.3004\) and
\(0.01002\) at the endpoint. The two bounds follow by recurrence.

For arbitrary \(a\), a single prime \(p\equiv a\bmod q\) can dominate the Riesz
main term \(y^2/2\) through \(h\log p\,(y-p)\) only if
\[
  p<\tfrac12 h\log p,\qquad
  \max_{y\ge p}\frac{y-p}{y^2}=\frac1{4p}.
\]
The equality follows from \((y-2p)^2\ge0\), with equality at \(y=2p\). By
monotonicity of \(p/\log p\) for \(p>e\), this restricts
\(p\le(1/2+o(1))q\log q\). The prime number theorem bounds the number of classes
such primes can represent by \((1/2+o(1))(q-1)\).
The full sum in (3) and the inverse-class sum in (4) retain the contributions
discarded by this test. A proof for all classes still requires both profile
signs uniformly in the remaining progressions, or another mechanism.

\newpage
\head{Checks at the first prime beyond the published range}
The next prime after the previous range is \(10\,009\). At this modulus the two
arithmetic identities apply to 10,006 of the 10,008 reduced classes. A scan of
65 integer scales \(1\le y\le2\,000\,000\), using both signs of \(\log y\),
gave the following strict inequalities for the profile.
\begin{center}
\begin{tabular}{lrrr}
\hline
Evaluation & Negative & Positive & Both signs\\
\hline
Arithmetic identities (3) and (4) & 10,006 & 9,779 & 9,779\\
One-ordinate certificate, all classes & 10,008 & 10,008 & 10,008\\
\hline
\end{tabular}
\end{center}
The first row used 128-bit and 160-bit outward-rounded intervals. The first
recorded witness for each sign agreed at both precisions. In the 227 classes
where this grid gave only a negative value, 98 were squares and 129 were nonsquares.
As a separate test of (3)--(4), the Hurwitz-transform constants at 1 and 2 agreed
with 20 direct logarithmic-derivative calculations at modulus 11. All 160 profile
evaluations for the eight applicable classes agreed with a sum over 850 certified
positive ordinates to within the published absolute tail bound \(0.154\).

For the second row the positive ordinate of Conrey character 9959 lies in
\((402441651\pm2)/2^{44}\). The two published per-prime programs verified
opposite endpoint signs, exclusion of all other inducing characters at the
ordinate, and absence of a family zero at twice its height. The dimension-three
kernel has amplitude \(>1.8856\). With
\href{https://arxiv.org/pdf/1305.3087#page=14}{Platt's Theorems 7.1 and 7.2}
and the dichotomy in the previous paper, the bounds
\(\liminf E_{q,a}<-3/2\), \(\limsup E_{q,a}>1/4\) hold unconditionally for all
10,008 classes at this prime. This one-modulus verification is separate from
the argument for unbounded primes.

The endpoint inequalities for the explicit class and the powers of two were
checked at 160 bits. Direct arithmetic evaluations gave the stated margins for
\(a=2,8\) at \(q=10\,009\) and \(a=4,16\) at \(q=10\,037\).
Lean 4.33.1 separately checks the two elementary prime-power exclusions and the
exact threshold for one prime to exceed the Riesz main term.

\textbf{Prior work.} Johnston posed the classification in
\href{https://arxiv.org/pdf/2407.03530#page=9}{Problem 23}.
Zhao treats the same integral in \href{https://arxiv.org/html/2411.18903v2#S8}{equation (35)}.
Erd\H{o}s's positive-proportion theorem concerns least primes.
\href{https://arxiv.org/abs/1710.00088}{Martin and Ng} study ordinary prime races
under zero-independence conditions.
\href{https://arxiv.org/abs/2010.05944}{De la Bretèche and Fiorilli} obtain
results on moments over residue classes for weighted prime sums. A search
covering those works and other smoothed prime races located no prior proof of
(1) or the rule selecting class 2 or 4 under Dirichlet GRH alone.
Priority remains unconfirmed.
\end{document}
